% Generated draft: requires mathematical and source review.
\documentclass[11pt]{article}
\usepackage[margin=25mm]{geometry}
\usepackage{fontspec}
\setmainfont{Latin Modern Roman}
\usepackage{amsmath,amssymb,mathtools,unicode-math}
\setmathfont{Latin Modern Math}
\usepackage{enumitem,needspace,xurl,hyperref}
\setlength{\parindent}{0pt}
\setlength{\parskip}{5pt}
\setlength{\emergencystretch}{3em}
\newcommand{\N}{\mathbb{N}}
\newcommand{\Z}{\mathbb{Z}}
\newcommand{\Q}{\mathbb{Q}}
\newcommand{\R}{\mathbb{R}}
\newcommand{\C}{\mathbb{C}}
\newcommand{\F}{\mathbb{F}}
\newcommand{\A}{\mathbb{A}}
\newcommand{\PP}{\mathbb{P}}
\newcommand{\E}{\mathbb{E}}
\newcommand{\eps}{\varepsilon}
\newcommand{\dd}{\,\mathrm{d}}
\begin{document}
\section*{UnsolvedMath ID: 30005776}
\textbf{Source problem code:} catalogue entry
\subsection{Definitions and mathematical statement}
Work over \(\C\). For a finite group \(G\), a left-left Yetter--Drinfeld module \(V\in{}_{\C G}^{\C G}\mathcal{YD}\) is a complex vector space with a left \(G\)-action and a grading
\[
V=\bigoplus_{g\in G}V_g,
\qquad h\cdot V_g\subseteq V_{hgh^{-1}}\quad(g,h\in G).
\]
Its braiding is
\[
c(v\otimes w)=g\cdot w\otimes v
\qquad(v\in V_g,\ w\in V).
\]
Give the tensor algebra \(T(V)=\bigoplus_{n\geq0}V^{\otimes n}\) its graded braided Hopf algebra structure in which every element of \(V\) is primitive. Let \(I(V)\) be the largest homogeneous braided Hopf ideal of \(T(V)\) contained in \(\bigoplus_{n\geq2}V^{\otimes n}\). The Nichols algebra of \(V\) is
\[
\mathcal B(V)=T(V)/I(V).
\]
Thus \(\mathcal B(V)_1\cong V\); in particular, finite dimensionality of \(\mathcal B(V)\) forces finite dimensionality of \(V\).

\textbf{Classification and relations problem.} Classify all pairs \((G,V)\) such that \(G\) is a finite group, \(V\in{}_{\C G}^{\C G}\mathcal{YD}\), and
\[
\dim_{\C}\mathcal B(V)<\infty.
\]
For every pair in this classification, determine explicit defining relations: exhibit a set \(R\subseteq I(V)\) generating \(I(V)\) as a two-sided ideal, so that
\[
\mathcal B(V)\cong T(V)/\langle R\rangle.
\]
No restriction to abelian groups, simple groups, irreducible modules, faithful actions, or nonzero modules is stated in the supplied question. The source does not specify the equivalence relation intended by ``classify.'' In particular, it does not say whether groups are fixed or identified by group isomorphisms together with compatible transport of the action and grading; this convention remains to be clarified.

\textbf{Supplied contextual qualifications.} The report excerpt attributes the classification for abelian \(G\) to Heckenberger and the corresponding defining relations to Angiono. These are reported background claims, not independently verified here. It suggests finite simple groups and modules with simple rack support as directions of study. Here \(\operatorname{supp}V=\{g\in G:V_g\ne0\}\) carries the conjugation rack operation \(x\triangleright y=xyx^{-1}\). The excerpt describes simple racks as those of prime-power cardinality \(p^n\), with \(p\) prime, and the remaining simple racks, including conjugacy classes in finite simple groups. It reports a theorem labeled [9]: if \(\dim\mathcal B(V)<\infty\) and \(\operatorname{supp}V\) is a simple rack of prime cardinality \(p\), then \(V\) lies on a short list and \(p\leq7\). Neither that list nor the bibliographic identification of [9] is supplied. Its further suggestion that the techniques extend to simple racks of cardinality \(p^n\) is explicitly based on unpublished calculations and is not asserted here as a theorem. These contextual passages do not restrict the classification problem above.
\subsection{Short English statement}
Classify, over \(\C\), all Yetter--Drinfeld modules over all finite groups whose Nichols algebras are finite dimensional, and give defining relations for each algebra. The intended equivalence convention for the classification is unspecified. The supplied record labels the general problem open and reports the abelian case as solved; these are supplied assessments requiring independent review.
\subsection{Sources}
\begin{itemize}
\item Local registry: \texttt{data/unsolvedmath/problems.json}, with \texttt{catalog/problems.json} as the additional-record registry; supplied complete record for permanent UnsolvedMath ID \(30005776\), titled ``Finite-Dimensional Nichols Algebras over Finite Groups,'' exact source code \texttt{catalogue entry}. Registry files were not independently inspected for this draft.
\item Supplied primary citation: \emph{Mini-Workshop: Bridging Number Theory and Nichols Algebras via Deformations} (2024), Oberwolfach Reports, pp.\ 235--272, DOI \texttt{10.4171/owr/2024/5}, \url{https://doi.org/10.4171/owr/2024/5}. The supplied excerpt identifies the target as Question 3 in the discussion headed ``2. Racks.'' No individual proposer is supplied.
\item Supplied literature source labeled ``original OWR report'': \url{https://doi.org/10.4171/owr/2024/5}.
\item Supplied literature source labeled ``Heckenberger--Vendramin, classification over non-abelian groups'': \url{https://arxiv.org/abs/1412.0857}.
\item Literature assessment date: \texttt{2026-08-22}. This is supplied dated curation, not an independently verified literature assessment. The supplied assessment labels the general classification and relations problem open and says that substantial rank-specific and rack-specific classifications do not settle it.
\item The record labels its statement \texttt{corrected\_verified} and says that it was reconstructed and checked against the official Oberwolfach report. These are supplied provenance claims; no source visit or independent verification is claimed here. The excerpt attributes the abelian classification to Heckenberger and its relations to Angiono, and refers to talks by Angiono and Vendramin without supplying their individual citations.
\end{itemize}
\end{document}
